\section{Introduction}

When an agent ends a turn without building, did it preserve cash or fail to execute a purchase? When it declares bankruptcy, was payment impossible or merely absent from the current menu? When a rival accepts its proposal, did the rival believe the proposal's public premise? These questions concern different properties of the same trajectory. The applied action alone need not settle any of them. An incorrect behavioral label misdescribes the current decision, while the resolved action changes the opportunities and information available later.

Economic environments make these questions concrete. Vending-Bench evaluates business operation, and its Arena communication study checks statements against simulator records and analyzes promise follow-through~\cite{backlund2025vendingbench,li2026arena}. AgentMonopoly and FinBoardBench study bargaining and resource use in board games; DECE compares complete model--harness configurations across Eurogames~\cite{wu2026agentmonopoly,hu2026finboardbench,zhang2026dece}. These studies already motivate process-sensitive evaluation and state-grounded communication analysis. We address a narrower problem: which evidence is sufficient for a particular claim about an observed decision, and how can that sufficiency be checked?

MonopolyBench records legal menus, every response attempt, resolution conditions, delivered memory and economic events in a deterministic four-player environment. We organize the audit around three distinctions. \emph{Resolution provenance} separates submitted responses from engine defaults and conflicting received channels. \emph{Reachability} separates a current menu or nominal balance sheet from a legal sequence ending at a stated goal. \emph{Delivery and response} separate logged speech, recipient exposure, textual endorsement and an identified effect on choice. These are evidentiary distinctions, not interchangeable behavioral scores.

We analyze twenty completed games with 6,989 decisions and 8,679 attempts. The most direct test is a retrospective census of all sixty eliminations and sixty-six engine defaults. Each claim specifies its historical prefix, intervention class and endpoint. A successful own-asset liquidation witnesses current-debt payment. A generous asset bound below that debt proves impossibility within the same class. Failed search supplies no such proof unless it exhausts the permitted continuations. A portable suite replays all original events and verifies these certificates.

The distinctions change what the records establish. Sixty-three defaulted construction decisions contain a reachable target, while the submitted encoding remains invalid. Three eliminations admit payment; one has enough nominal equity but no legal liquidation route. Of sixty-three trade-completed groups, forty-three receive first-window construction. The twenty remaining groups include nine without a build opportunity, four ending in defaults, four in content/tool conflicts and three in valid END records. Aggregating them as investment restraint would erase materially different events. These are cohort-specific counts, not error rates in other benchmarks.

The social analysis establishes a different kind of result. A recipient who explicitly corrects a false premise and accepts the transaction is a counterexample to treating settlement as sufficient evidence of endorsement. Other cases remain unidentified: a buyer echoes inheritance framing, but the record does not show that it believes a false financial minimum or would have paid less without the message. Engine branches test the financial route; recipient prompts and responses constrain its interpretation. Neither reveals latent intent.

The contribution is a claim-oriented audit procedure, an executable finite-corpus demonstration, and source-linked social counterexamples that mark its limits. It goes beyond replaying recorded actions by specifying what an alternative sequence proves. It goes beyond reading explanations by checking the action and information delivered. The engine and saved games are publicly available, and this revision supplies a self-contained certificate suite tested outside the repository. Economic trajectories explain why the distinctions matter. We make no intrinsic model ranking, general deception rate or capability--alignment trend claim: public deception was permitted, and capability and response policies were not independently varied.

